\clearpage\corrigesChapitre{20}{Capteurs et outils numériques}
\corrigesSouples
\begin{blocLivre}{cor-20-1}\begin{corrige}{1}{Reconnaître un capteur}
\begin{questions}
\item Un capteur peut détecter l’éclairement ; un autre la température.
\item Non : sa réponse physique est convertie, acquise et traitée avant affichage.
\end{questions}
\end{corrige}\end{blocLivre}
\begin{blocLivre}{cor-20-2}\begin{corrige}{2}{Lire une courbe}
\schema[corrige]{capteur-courbe}
\begin{questions}
\item $R=20,0-0,400\times20=\qty{12,0}{k\ohm}$.
\item La résistance diminue lorsque la température augmente sur le domaine étudié.
\end{questions}
\end{corrige}\end{blocLivre}
\begin{blocLivre}{cor-20-3}\begin{corrige}{3}{Une température inconnue}
\begin{questions}
\item $\theta=(20,0-10,0)/0,400=25,0\,^{\circ}\mathrm C$.
\item 25 °C est entre 0 et 40 °C : on reste dans le domaine d’étalonnage.
\end{questions}
\end{corrige}\end{blocLivre}
\begin{blocLivre}{cor-20-4}\begin{corrige}{4}{Une conversion numérique}
\begin{questions}
\item $U=5,00\times600/1000=\qty{3,00}{V}$.
\item $k=1000\times1,25/5,00=250$.
\end{questions}
\end{corrige}\end{blocLivre}
\begin{blocLivre}{cor-20-5}\begin{corrige}{5}{Un pont simple}
\schema[corrige]{pont-capteur}
\begin{questions}
\item $I=5,00/20000=\qty{0,250}{mA}$.
\item $U=R_cI=10000\times0,000250=\qty{2,50}{V}$.
\end{questions}
\end{corrige}\end{blocLivre}
\begin{blocLivre}{cor-20-6}\begin{corrige}{6}{La grandeur utile}
\begin{questions}
\item L’éclairement reçu par le capteur.
\item Il faut connaître la correspondance entre réponse électrique et grandeur recherchée, dans les conditions d’utilisation.
\end{questions}
\end{corrige}\end{blocLivre}
\begin{blocLivre}{cor-20-7}\begin{corrige}{7}{Une résistance retrouvée}
\schema[corrige]{pont-capteur}
\begin{questions}
\item $R_c=10,0\times2,00/(5,00-2,00)=\qty{6,67}{k\ohm}$.
\item $\theta=(20,0-6,67)/0,400\simeq33,3\,^{\circ}\mathrm C$.
\end{questions}
\end{corrige}\end{blocLivre}
\begin{blocLivre}{cor-20-8}\begin{corrige}{8}{Hors domaine}
\begin{questions}
\item Le modèle a seulement été testé de 0 à 40 °C ; le prolongement est une extrapolation.
\item Vérifier les limites du composant et effectuer un étalonnage adapté sur le nouveau domaine avec une référence.
\end{questions}
\end{corrige}\end{blocLivre}
\begin{blocLivre}{cor-20-9}\begin{corrige}{9}{Un seuil à programmer}
\begin{questions}
\item $R_s=20,0-0,400\times30,0=\qty{8,00}{k\ohm}$.
\item Pour $R<8,00$ kΩ, puisque la résistance diminue quand la température augmente.
\end{questions}
\end{corrige}\end{blocLivre}
\begin{blocLivre}{cor-20-10}\begin{corrige}{10}{Choisir les axes en Python}
\begin{questions}
\item La température en abscisse, la résistance en ordonnée.
\item Les noms des grandeurs, les unités, les points mesurés, la légende du modèle et le domaine de validité.
\end{questions}
\end{corrige}\end{blocLivre}
\begin{blocLivre}{cor-20-11}\begin{corrige}{11}{Une division impossible}
\begin{questions}
\item Le dénominateur est nul : la formule ne donne aucune résistance finie.
\item Signaler une mesure hors domaine ou un défaut possible, vérifier le circuit et éviter la division. Il ne faut pas afficher une température arbitraire.
\end{questions}
\end{corrige}\end{blocLivre}
\begin{blocLivre}{cor-20-12}\begin{corrige}{12}{Deux sources d’erreur}
\demarche{Vérifier successivement le capteur, les connexions et l’acquisition.}
\begin{questions}
\item Mauvais contact, fluctuations de l’alimentation, bruit électrique ou stabilisation thermique insuffisante.
\item Vérifier les connexions et la référence de tension, attendre l’équilibre puis répéter avec un thermomètre de référence.
\end{questions}
\retiens{Une valeur instable peut venir de la chaîne de mesure et pas seulement de la grandeur étudiée.}
\end{corrige}\end{blocLivre}

